\section{Conclusion}
\label{sec:clustering_agent_representations_conclusion}

The \gls{car} method was developed to investigate how concepts are organized within the representations of \gls{rl} policies.
It is based on three hypotheses introduced throughout this thesis, the \MechanisticHypothesisName, the \ConceptHypothesisName, and the \RepresentationConvergenceHypothesisName.

\gls{car} translates these assumptions into an empirical analysis.
Multiple surrogate policies are trained independently to reproduce the same policy.
Their \glspl{lr} are then clustered and compared to assess whether similar conceptual structures emerge across policies.
The \ClusteringUniversalityName and \AbstractionScoreName provide quantitative measures for evaluating the consistency and abstraction of the resulting clusters.

This makes \gls{car} a general procedure for analyzing agent representations.
The analysis and evaluation metrics are fully automated and require no manual intervention or environment-specific adjustments.
The method can therefore be applied to a wide range of \gls{rl} settings and provides a basis for systematically studying the internal representations of neural policies.

The next chapter applies \gls{car} to several environments to evaluate its ability to identify meaningful conceptual structures.

